\documentclass{article}
\usepackage[english]{babel}
\usepackage{amsmath}
\usepackage{amsthm}
\usepackage{amssymb}

\newtheorem{theorem}{Theorem}[section]
\newtheorem{lemma}[theorem]{Lemma}
\newenvironment{solution}{
\begin{proof}[Solution]}{
\end{proof}}

\begin{document}

% ============================================================
\section*{Section 1.1}

\subsection*{1.1.8}
Write the number in the form $a + bi$:
\[
  \frac{(8+2i)-(1-i)}{(2+i)^2}
\]
\begin{solution}
  Rearranging:
  \begin{align*}
    \frac{(8+2i)-(1-i)}{(2+i)^2}
      &= \frac{(8+2i)-(1-i)}{3+4i} \\
      &= \frac{(8+2i)-(1-i)}{3+4i} \cdot \frac{3-4i}{3-4i} \\
      &= \frac{33-19i}{25} \\
      &= \frac{33}{25} - \frac{19}{25}i
  \end{align*}
\end{solution}

\subsection*{1.1.10}
Write the number in the form $a + bi$:
\[
  \left[\frac{2+i}{6i-(1-2i)}\right]^2
\]
\begin{solution}
  \begin{align*}
    \left(\frac{2+i}{6i-(1-2i)}\right)^2
      &= \frac{(2+i)^2}{(-1+8i)^2} \\
      &= \frac{4+4i-1}{1-16i-64} \\
      &= \frac{3+4i}{-63-16i} \\
      &= \frac{(3+4i)(-63+16i)}{(-63-16i)(-63+16i)} \\
      &= \frac{-189+48i-252i-64}{63^2+16^2} \\
      &= \frac{-253-204i}{63^2+16^2} \\
      &= \frac{-253}{63^2+16^2} + \frac{-204i}{63^2+16^2}
  \end{align*}
\end{solution}

\subsection*{1.1.14}
Show that $\operatorname{Re}(iz) = -\operatorname{Im} z$ for every complex number $z$.
\begin{proof}
  Write $z = x + yi$ where $\operatorname{Re}(z) = x$ and $\operatorname{Im}(z) = y$.
  Then $iz = xi - y$ and $\operatorname{Re}(iz) = -y = -\operatorname{Im}(z)$.
\end{proof}

\subsection*{1.1.15}
Let $k$ be an integer. Show that
\[
  i^{4k} = 1, \quad i^{4k+1} = i, \quad i^{4k+2} = -1, \quad i^{4k+3} = -i.
\]
\begin{proof}
  \textbf{1.}
  \[
    i^{4k} = (i^2)^{2k} = ((-1)^2)^k = 1^k
  \]
  Because $1^k = 1 \quad \forall k \in \mathbb{Z}$.

  \textbf{2.}
  \begin{align*}
    i^{4k+1} &= i^{4k} \cdot i \\
      &= 1 \cdot i && \text{we can substitute the result of 1.1.15 part 1} \\
      &= i
  \end{align*}

  \textbf{3.}
  \begin{align*}
    i^{4k+2} &= i^{4k} \cdot i^2 \\
      &= 1 \cdot i^2 && \text{using the result from 1.1.15 part 1} \\
      &= i^2 \\
      &= -1
  \end{align*}

  \textbf{4.}
  \begin{align*}
    i^{4k+3} &= i^{4k} \cdot i^3 \\
      &= 1 \cdot i^3 && \text{using the result from 1.1.15 part 1} \\
      &= i^3 \\
      &= i^2 \cdot i \\
      &= -1 \cdot i \\
      &= -i
  \end{align*}
\end{proof}

\subsection*{1.1.20}
Solve each of the following equations for $z$.

\subsubsection*{1.1.20 (c)}
\[
  (2-i)z + 8z^2 = 0
\]
\begin{solution}
  \begin{align*}
    (2-i)z + 8z^2 &= 0 \\
    z\bigl((2-i) + 8z\bigr) &= 0
  \end{align*}
  Thus:
  \[
    z = 0 \text{ or } z = \frac{-2+i}{8}
  \]
\end{solution}

\subsection*{1.1.22}
Find all solutions to the equation $z^4 - 16 = 0$.
\begin{solution}
  \begin{align*}
    z^4 - 16 &= 0 \\
    (z^2-4)(z^2+4) &= 0
  \end{align*}
  Thus we solve each root independently:
  \begin{gather*}
    (z^2-4) = 0 \implies z = \pm 2 \\
    (z^2+4) = 0 \implies z = \pm 2i
  \end{gather*}
  Thus all solutions are: $z = \{2, -2, 2i, -2i\}$.
\end{solution}

\subsection*{1.1.24}
Let $z$ be a complex number such that $\operatorname{Im} z > 0$. Prove that $\operatorname{Im}(1/z) < 0$.
\begin{proof}
  Let $z = x + yi$ with $\operatorname{Im}(z) > 0$, so $y > 0$. Thus:
  \begin{align*}
    \frac{1}{z} &= \frac{1}{x+yi} \\
      &= \frac{x-yi}{x^2+y^2} && \text{multiplying by the conjugate} \\
      &= \frac{x}{x^2+y^2} + \frac{-yi}{x^2+y^2} && \text{separating the real and imaginary components}
  \end{align*}
  Thus taking the imaginary component:
  \[
    \operatorname{Im}\left(\frac{1}{z}\right) = \frac{-y}{x^2+y^2}
  \]
  The denominator is positive because any $a^2$ for $a \in \mathbb{R}$ is positive and sums of positives are also positive.
  The numerator is negative because $y > 0$ so $-y < 0$.
  Finally we know the whole fraction is negative: $u, v \in \mathbb{R},\ u < 0,\ v > 0 \implies \frac{u}{v} < 0$.
  Thus:
  \[
    \operatorname{Im}\left(\frac{1}{z}\right) < 0
  \]
\end{proof}

% ============================================================
\section*{Section 1.2}

\subsection*{1.2.7}
Describe the set of points $z$ in the complex plane that satisfies each of the following.

\subsubsection*{1.2.7 (d)}
\[
  |z-1| = |z+i|
\]
\begin{solution}
  \begin{align*}
    |z-1| &= |z+i| \\
    \sqrt{(x-1)^2 + y^2} &= \sqrt{x^2 + (y+1)^2} && \text{substituting } z = x+yi \\
    (x-1)^2 + y^2 &= x^2 + (y+1)^2 && \text{squaring both sides} \\
    x^2 - 2x + 1 + y^2 &= x^2 + y^2 + 2y + 1 \\
    -2x &= 2y && \text{canceling terms} \\
    y &= -x
  \end{align*}
  Thus the set satisfying $|z-1| = |z+i|$ is $\{x + yi : y = -x,\ x, y \in \mathbb{R}\}$, the perpendicular bisector of $1$ and $-i$.
\end{solution}

\subsubsection*{1.2.7 (e)}
\[
  |z| = \operatorname{Re} z + 2
\]
\begin{solution}
  \begin{align*}
    |z| &= \operatorname{Re}(z) + 2 \\
    \sqrt{x^2+y^2} &= x + 2 && \text{substituting } z = x+yi \text{ such that } \operatorname{Re}(z) = x \\
    x^2 + y^2 &= x^2 + 4x + 4 && \text{squaring both sides} \\
    y^2 &= 4x + 4 \\
    y &= \pm\sqrt{4x+4}
  \end{align*}
  Thus the set satisfying $|z| = \operatorname{Re}(z) + 2$ is $\{x + yi : y = \pm\sqrt{4x+4},\ x, y \in \mathbb{R}\}$.
\end{solution}

\subsubsection*{1.2.7 (f)}
\[
  |z-1| + |z+1| = 7
\]
\begin{solution}
  We begin with:
  \[
    |z-1| + |z+1| = 7
  \]
  Substituting in the following for the absolutes and naming them $A$ and $B$ respectively:
  \begin{align*}
    A &= |z-1| = \sqrt{(x-1)^2 + y^2} \\
    B &= |z+1| = \sqrt{(x+1)^2 + y^2}
  \end{align*}
  We get:
  \begin{align*}
    A &= 7 - B \\
    A^2 &= (7-B)^2 && \text{we can square both sides} \\
    \intertext{and substituting the expression for $A$ from the definition of the absolute:}
    (x-1)^2 + y^2 &= (7-B)^2 \\
    (x-1)^2 + y^2 &= 49 - 14B + B^2 \\
    \intertext{substituting the expression for $B$ from the definition of the absolute:}
      &= 49 - 14B + (x+1)^2 + y^2 \\
    (x-1)^2 &= 49 - 14B + (x+1)^2 && \text{and canceling terms} \\
    x^2 - 2x + 1 &= 49 - 14B + x^2 + 2x + 1 \\
    14B &= 4x + 49 && \text{simplifying} \\
    \intertext{finally substituting in $B$ for its definition from the absolute and simplifying further:}
    49 &= 14\sqrt{(x+1)^2 + y^2} - 4x \\
    49^2 &= \left(14\sqrt{(x+1)^2 + y^2} - 4x\right)^2 \\
    49^2 + 392x + 16x^2 &= 14^2(x^2 + 2x + 1 + y^2) \\
    0 &= 180x^2 + 196 - 49^2 + 169y^2 \\
    2205 &= 180x^2 + 196y^2
  \end{align*}
  Thus the set satisfying $|z-1| + |z+1| = 7$ is:
  \[
    \{x + yi : 2205 = 180x^2 + 196y^2,\ x, y \in \mathbb{R}\}
  \]
\end{solution}

\subsection*{1.2.16}
Prove that if $|z| = 1$ $(z \neq 1)$, then
\[
  \operatorname{Re}\left[\frac{1}{1-z}\right] = \frac{1}{2}.
\]
\begin{proof}
  Rearranging $\frac{1}{1-z}$ to separate the real and imaginary components, we get:
  \begin{align*}
    \frac{1}{1-z} &= \frac{1}{1-(x+yi)} \\
      &= \frac{1}{(1-x) - yi} \cdot \frac{(1-x) + yi}{(1-x) + yi} \\
    \intertext{which is well defined because $z \neq 1 \implies$ the denominator $(1-x) + yi \neq 0$}
      &= \frac{(1-x) + yi}{(1-x)^2 + y^2} \\
      &= \frac{1-x}{(1-x)^2 + y^2} + \frac{yi}{(1-x)^2 + y^2}
  \end{align*}
  Thus taking just the real component:
  \[
    \operatorname{Re}\left(\frac{1}{1-z}\right) = \frac{1-x}{(1-x)^2 + y^2}
  \]
  Now from $|z| = 1$ we can expand the absolute and rearrange algebraically:
  \begin{align*}
    1 &= |z| \\
      &= \sqrt{x^2+y^2} \\
      &= x^2 + y^2 \\
    2 - 2x &= x^2 - 2x + 1 + y^2 \\
      &= (x-1)^2 + y^2 \\
    1 &= \frac{2-2x}{(x-1)^2 + y^2} \\
    \intertext{denominator is nonzero because $z \neq 1 \implies (x-1)^2 + y^2 \neq 0$}
    \frac{1}{2} &= \frac{1-x}{(x-1)^2 + y^2}
  \end{align*}
  Which is exactly $\operatorname{Re}\left(\frac{1}{1-z}\right)$. Thus:
  \[
    \operatorname{Re}\left(\frac{1}{1-z}\right) = \frac{1}{2}
  \]
\end{proof}

% ============================================================
\section*{Section 1.3}

\subsection*{1.3.7}
Find the argument of each of the following complex numbers and write each in polar form.

\subsubsection*{1.3.7 (d)}
\[
  -2\sqrt{3} - 2i
\]
\begin{solution}
  $z = -2\sqrt{3} - 2i$.
  Find the argument angle $\theta$ and account for the fact that arctan's range is always within the 1st or 4th quadrants:
  \begin{gather*}
    \tan\theta = \frac{-2}{-2\sqrt{3}} \implies \theta = \frac{\pi}{6} \\
    \theta + \pi = \frac{7\pi}{6} \text{ or equivalently } -\frac{5\pi}{6}
  \end{gather*}
  The absolute is:
  \[
    |-2\sqrt{3} - 2i| = \sqrt{(-2\sqrt{3})^2 + (-2)^2} = 4
  \]
  Thus we write:
  \[
    z = 4\left(\cos\left(-\frac{5\pi}{6}\right) + i\sin\left(-\frac{5\pi}{6}\right)\right)
  \]
\end{solution}

\subsubsection*{1.3.7 (g)}
\[
  \frac{-1+\sqrt{3}\,i}{2+2i}
\]
\begin{solution}
  We begin by separating the real and imaginary components of $z = \frac{-1+\sqrt{3}i}{2+2i}$ by multiplying the numerator and denominator by the conjugate of the denominator:
  \[
    z = \frac{(-2 + 2\sqrt{3}) + i(2 + 2\sqrt{3})}{8}
  \]
  Then taking the inverse tangent:
  \[
    \theta = \arctan\left(\frac{2+2\sqrt{3}}{-2+2\sqrt{3}}\right) = \frac{5}{12}\pi \quad \text{which is already in the 1st quadrant}
  \]
  And computing the absolute:
  \[
    r = |z| = \frac{\sqrt{2}}{2}
  \]
  Finally:
  \[
    z = \frac{\sqrt{2}}{2}\left(\cos\left(\frac{5}{12}\pi\right) + i\sin\left(\frac{5}{12}\pi\right)\right)
  \]
\end{solution}

\subsection*{1.3.10}
Show the following:

\subsubsection*{1.3.10 (a)}
\[
  \arg(z_1 z_2 z_3) = \arg z_1 + \arg z_2 + \arg z_3
\]
\begin{proof}
  We know that for two complex numbers $z_1 = r_1(\cos\theta_1 + i\sin\theta_1),\ z_2 = r_2(\cos\theta_2 + i\sin\theta_2) \in \mathbb{C}$:
  \[
    z_1 \cdot z_2 = r_1 r_2[\cos(\theta_1 + \theta_2) + i\sin(\theta_1 + \theta_2)]
  \]
  Thus, we can apply this to $z_1, z_2, z_3$ by writing $z_1 \cdot z_2$ as $z_a$:
  \[
    z_a = r_a(\cos(\theta_a) + i\sin(\theta_a)) \quad \text{with } r_a = r_1 \cdot r_2, \quad \theta_a = \theta_1 + \theta_2
  \]
  Then:
  \begin{align*}
    z_1 z_2 z_3 = z_a \cdot z_3
      &= r_a r_3[\cos(\theta_a + \theta_3) + i\sin(\theta_a + \theta_3)] \\
      &= r_1 r_2 r_3[\cos(\theta_1 + \theta_2 + \theta_3) + i\sin(\theta_1 + \theta_2 + \theta_3)]
  \end{align*}
  Taking the argument from the $\sin$ and $\cos$ terms:
  \[
    \arg(z_1 z_2 z_3) = \theta_1 + \theta_2 + \theta_3
  \]
\end{proof}

\subsubsection*{1.3.10 (b)}
\[
  \arg\left(\frac{z_1}{z_2}\right) = \arg z_1 - \arg z_2
\]
\begin{proof}
  We can write the argument of $z_2 = x + iy \in \mathbb{C}$ as $\arg(z_2) = \arctan\left(\frac{y}{x}\right) = \theta_2$.

  Taking the conjugate of $z_2$, then solving for its argument:
  \[
    \overline{z_2} = x + i(-y)
  \]
  Which can be written as $r_2(\cos\theta + i\sin\theta)$ for some angle $\theta$. Solving for $\theta$ by setting the real and imaginary parts equal:
  \begin{align*}
    x &= r_2\cos\phi \\
    y &= r_2\sin\phi \\
    \frac{-y}{x} &= \tan\phi
  \end{align*}
  This leaves:
  \begin{gather*}
    \theta = \arctan\left(-\frac{y}{x}\right) \implies -\phi = \arctan\left(\frac{y}{x}\right) = \theta_2 \\[1ex]
    \therefore \quad z = r_2(\cos(-\theta_2) + i\sin(-\theta_2))
  \end{gather*}
  Now we can apply the known identity $z_1 \cdot z_2 = r_1 r_2[\cos(\theta_1 + \theta_2) + i\sin(\theta_1 + \theta_2)]$ such that:
  \[
    z_1\overline{z_2} = r_1 r_2[\cos(\theta_1 - \theta_2) + i\sin(\theta_1 - \theta_2)]
  \]
  and we can read off the argument:
  \[
    \arg(z_1\overline{z_2}) = \theta_1 - \theta_2 = \arg(z_1) - \arg(z_2)
  \]
\end{proof}

\subsection*{1.3.12}
Find the following.

\subsubsection*{1.3.12 (b)}
\[
  \operatorname{Arg}(-\pi)
\]
\begin{solution}
  $-\pi$ is purely real valued and negative so its argument is $\arg_0 = \pi$.
\end{solution}

\subsubsection*{1.3.12 (c)}
\[
  \operatorname{Arg}(10i)
\]
\begin{solution}
  $10i$ is purely imaginary so its argument is $\arg_0 = \frac{\pi}{2}$.
\end{solution}

% ============================================================
\section*{Additional Problem}

\subsection*{$\operatorname{Re}(zw) = \operatorname{Re}(z) \cdot \operatorname{Re}(w)$}
\begin{proof}
  It is not true that $\operatorname{Re}(zw) = \operatorname{Re}(z)\operatorname{Re}(w)$.

  Counterexample: $z = 1 + i, \quad w = 1 + i$.

  Plugging into the LHS:
  \[
    \operatorname{Re}(zw) = \operatorname{Re}((1+i)(1+i)) = \operatorname{Re}(1 + 2i - 1) = \operatorname{Re}(2i) = 0
  \]
  Plugging into the RHS:
  \[
    \operatorname{Re}(1+i) \cdot \operatorname{Re}(1+i) = 1
  \]
  Thus: $\operatorname{Re}(zw) \neq \operatorname{Re}(z)\operatorname{Re}(w)$.
\end{proof}

\end{document}